% GENERATED FILE. Edit canonical lessons, then run npm run generate:books.
% canonical-source-hash: fnv1a64:a5accfdcf20f0635
% canonical-lessons: PT-C401-faixa, PT-C401-berma, PT-C401-separador, PT-C401-no, PT-C401-acelerar

\chapter{Na autoestrada}
\label{ch:pt-na-autoestrada}
\hlchaptermodality{401}

\begin{chapteropening}
\textbf{By the end of this chapter:} \emph{I can talk about a motorway trip: the lanes, the hard shoulder, the central reservation, the junction and speeding up.}

Complete the last lesson of chapter 401: I can talk about a motorway trip: the lanes, the hard shoulder, the central reservation, the junction and speeding up.
\end{chapteropening}

\section[a faixa]{a faixa — the lane (of a road); the strip, the band}

\label{lesson:PT-C401-faixa}



\begin{quote}
Before the new one: say the Portuguese for a jungle, then the Portuguese for a slipper.
\end{quote}

\subsection*{You'll want to know: a faixa}
\textbf{a faixa} — \textquotedblleft{}the lane (of a road); the strip, the band\textquotedblright{}.

\begin{quote}
Na autoestrada, os carros lentos vão na faixa da direita.
\end{quote}

\begin{quote}
\emph{On the motorway, slow cars keep to the right-hand lane.}
\end{quote}

\emph{A faixa de rodagem} — the carriageway. \emph{Uma faixa} is also a strip of cloth or a banner.

\begin{cousinweb}
\emph{faixa} is inherited from Latin \emph{fascia} 'band, strip of cloth' (EtymDB): a lane is a long strip of road.
\end{cousinweb}

\begin{grammarlens}[title={What you've built}]
One more word for telling, explaining and describing at B1.
\end{grammarlens}

\subsection*{Your turn}
\begin{itemize}
\raggedright
  \item \emph{Say it:} \emph{a faixa}
  \item \emph{Say it:} \emph{a faixa}, once more
  \item \emph{Say it:} a sentence of your own with \emph{a faixa}
  \item \emph{Recall:} say the Portuguese for a jungle, then the Portuguese for a slipper, then say \emph{a faixa} again
\end{itemize}

\begin{tcolorbox}[breakable,colback=teal!4,colframe=teal!35!black,title={Before you move on}]
What does a faixa mean? (\textquotedblleft{}The lane (of a road); the strip, the band\textquotedblright{}.) Say it once more.
\end{tcolorbox}
\section[a berma]{a berma — the hard shoulder, the verge}

\label{lesson:PT-C401-berma}



\begin{quote}
Before the new one: say the Portuguese for a slipper, then the Portuguese for the lane.
\end{quote}

\subsection*{You'll want to know: a berma}
\textbf{a berma} — \textquotedblleft{}the hard shoulder, the verge\textquotedblright{}.

\begin{quote}
O carro parou na berma com um problema no motor.
\end{quote}

\begin{quote}
\emph{The car stopped on the hard shoulder with an engine problem.}
\end{quote}

In Brazil: \emph{o acostamento}.

\begin{cousinweb}
\emph{berma} was borrowed from French \emph{berme}, which came from Dutch \emph{berm} 'bank, raised edge' (EtymDB).
\end{cousinweb}

\begin{grammarlens}[title={What you've built}]
One more word for telling, explaining and describing at B1.
\end{grammarlens}

\subsection*{Your turn}
\begin{itemize}
\raggedright
  \item \emph{Say it:} \emph{a berma}
  \item \emph{Say it:} \emph{a berma}, once more
  \item \emph{Say it:} a sentence of your own with \emph{a berma}
  \item \emph{Recall:} say the Portuguese for a slipper, then the Portuguese for the lane, then say \emph{a berma} again
\end{itemize}

\begin{tcolorbox}[breakable,colback=teal!4,colframe=teal!35!black,title={Before you move on}]
What does a berma mean? (\textquotedblleft{}The hard shoulder, the verge\textquotedblright{}.) Say it once more.
\end{tcolorbox}
\section[o separador]{o separador — the central reservation; the divider, the tab}

\label{lesson:PT-C401-separador}



\begin{quote}
Before the new one: say the Portuguese for the lane, then the Portuguese for the hard shoulder.
\end{quote}

\subsection*{You'll want to know: o separador}
\textbf{o separador} — \textquotedblleft{}the central reservation; the divider, the tab\textquotedblright{}.

\begin{quote}
Há flores no separador da autoestrada.
\end{quote}

\begin{quote}
\emph{There are flowers on the central reservation of the motorway.}
\end{quote}

\emph{O separador central} is the strip between the two directions. In a browser or a folder, \emph{um separador} is a tab or a divider.

\begin{cousinweb}
\emph{separador} is built on \emph{separar} 'to separate', a learned borrowing from Latin \emph{sēparāre} 'to put apart' (EtymDB).
\end{cousinweb}

\begin{grammarlens}[title={What you've built}]
One more word for telling, explaining and describing at B1.
\end{grammarlens}

\subsection*{Your turn}
\begin{itemize}
\raggedright
  \item \emph{Say it:} \emph{o separador}
  \item \emph{Say it:} \emph{o separador}, once more
  \item \emph{Say it:} a sentence of your own with \emph{o separador}
  \item \emph{Recall:} say the Portuguese for the lane, then the Portuguese for the hard shoulder, then say \emph{o separador} again
\end{itemize}

\begin{tcolorbox}[breakable,colback=teal!4,colframe=teal!35!black,title={Before you move on}]
What does o separador mean? (\textquotedblleft{}The central reservation; the divider, the tab\textquotedblright{}.) Say it once more.
\end{tcolorbox}
\section[o nó]{o nó — the junction (of motorways); the knot}

\label{lesson:PT-C401-no}



\begin{quote}
Before the new one: say the Portuguese for the hard shoulder, then the Portuguese for the central reservation.
\end{quote}

\subsection*{You'll want to know: o nó}
\textbf{o nó} — \textquotedblleft{}the junction (of motorways); the knot\textquotedblright{}.

\begin{quote}
Saímos da autoestrada no próximo nó.
\end{quote}

\begin{quote}
\emph{We leave the motorway at the next junction.}
\end{quote}

Also a knot in a rope or a tie: \emph{dar um nó}.

\begin{cousinweb}
\emph{nó} is inherited from Latin \emph{nōdus} 'knot' (EtymDB): a junction is where roads are tied together.
\end{cousinweb}

\begin{grammarlens}[title={What you've built}]
One more word for telling, explaining and describing at B1.
\end{grammarlens}

\subsection*{Your turn}
\begin{itemize}
\raggedright
  \item \emph{Say it:} \emph{o nó}
  \item \emph{Say it:} \emph{o nó}, once more
  \item \emph{Say it:} a sentence of your own with \emph{o nó}
  \item \emph{Recall:} say the Portuguese for the hard shoulder, then the Portuguese for the central reservation, then say \emph{o nó} again
\end{itemize}

\begin{tcolorbox}[breakable,colback=teal!4,colframe=teal!35!black,title={Before you move on}]
What does o nó mean? (\textquotedblleft{}The junction (of motorways); the knot\textquotedblright{}.) Say it once more.
\end{tcolorbox}
\section[acelerar]{acelerar — to speed up, to accelerate}

\label{lesson:PT-C401-acelerar}



\begin{quote}
Before the new one: say the Portuguese for the central reservation, then the Portuguese for the junction.
\end{quote}

\subsection*{You'll want to know: acelerar}
\textbf{acelerar} — \textquotedblleft{}to speed up, to accelerate\textquotedblright{}.

\begin{quote}
Ela acelera para ultrapassar o autocarro.
\end{quote}

\begin{quote}
\emph{She speeds up to overtake the bus.}
\end{quote}

\emph{O acelerador} — the accelerator.

\begin{cousinweb}
\emph{acelerar} is a learned borrowing from Latin \emph{accelerāre} 'to hasten' (EtymDB), built on \emph{celer} 'swift' (Lewis \& Short).
\end{cousinweb}

\begin{grammarlens}[title={What you've built}]
One more word for telling, explaining and describing at B1.
\end{grammarlens}

\subsection*{Your turn}
\begin{itemize}
\raggedright
  \item \emph{Say it:} \emph{acelerar}
  \item \emph{Say it:} \emph{acelerar}, once more
  \item \emph{Say it:} a sentence of your own with \emph{acelerar}
  \item \emph{Recall:} say the Portuguese for the central reservation, then the Portuguese for the junction, then say \emph{acelerar} again
\end{itemize}

\begin{tcolorbox}[breakable,colback=teal!4,colframe=teal!35!black,title={Before you move on}]
What does acelerar mean? (\textquotedblleft{}To speed up, to accelerate\textquotedblright{}.) Say it once more.
\end{tcolorbox}
